\documentclass[preprint,11pt]{elsarticle}

\usepackage{amsmath,amssymb}
\usepackage{graphicx}
\usepackage{booktabs}
\usepackage{siunitx}
\usepackage[colorlinks=true,linkcolor=blue,citecolor=blue,urlcolor=blue]{hyperref}

\graphicspath{{figs/}}

\journal{---}

\begin{document}

\begin{frontmatter}

\title{Black-oil simulation of geologic CO$_2$ sequestration in MRST:
theory, verification practice, and a GNU Octave reproduction study}

\author[1]{Adriano Cortes}
\address[1]{---}

\begin{abstract}
This report documents the theory and the computational practice behind a
full reproduction of the geologic carbon storage (GCS) examples of
Sal\'o-Salgado et al. (2024), performed with the MATLAB Reservoir
Simulation Toolbox (MRST) under GNU Octave. We review the black-oil
formulation used to model CO$_2$--brine systems, the generation of PVT
input data from a thermodynamic mixing model, the Killough--Land model for
relative permeability hysteresis, and a Fickian model for molecular
diffusion of dissolved CO$_2$. Each model is exercised on the published
example decks: a PVT table verification, the PUNQ-S3 hysteresis validation
over 500 years, and a FluidFlower-inspired convective mixing problem at
two grid resolutions. We further include a spill-point structural trapping
analysis from the \texttt{co2lab} module. All results reproduce the
published behavior: the fluid object matches the tabulated PVT data to
machine precision; hysteresis immobilizes the CO$_2$ plume by residual
trapping and suppresses the mobile gas cap; and grid-level numerical
diffusion masks physical diffusion at coarse resolution, as predicted by a
simple scaling argument. We close with a portability study that documents
the code-level differences between MATLAB and Octave, including a silent
grid-corruption defect traced to a malformed block comment, and we discuss
its implications for reproducibility of simulation software. An appendix
further audits MRST's well-reservoir coupling on the diffusion case's
unstructured mesh: we show it approximates the well cell as an
axis-aligned Cartesian block rather than using its true PEBI geometry, and
quantify the resulting gap against a grid-aware alternative well index.
The full reproduction is also released as four executable Jupyter
notebooks, one per case, reusing MRST's Octave kernel.
\end{abstract}

\begin{keyword}
geologic carbon storage \sep black-oil model \sep MRST \sep relative
permeability hysteresis \sep molecular diffusion \sep reproducibility
\end{keyword}

\end{frontmatter}

\section{Introduction}
\label{sec:intro}

Geologic carbon storage (GCS) injects supercritical CO$_2$ into deep
saline aquifers. Four mechanisms retain the injected CO$_2$: structural
and stratigraphic trapping under low-permeability seals, residual
trapping as disconnected ganglia after imbibition, solubility trapping by
dissolution into brine, and mineral trapping over long times
\cite{benson2005}. Numerical simulation is the primary tool to assess
these mechanisms at field scale. The MATLAB Reservoir Simulation Toolbox
(MRST) \cite{lie2019,lie-moyner2021} provides an open-source platform for
such simulations, with full access to the discretization, the nonlinear
solvers, and the automatic differentiation (AD) backends.

Sal\'o-Salgado et al.~\cite{salo2024} extended the \texttt{ad-blackoil}
module of MRST with three capabilities that increase the fidelity of
black-oil GCS models: (i) generation of CO$_2$--brine PVT input data from
a thermodynamic mixing model; (ii) relative permeability hysteresis of
the nonwetting phase, which controls residual trapping; and (iii)
molecular diffusion of dissolved CO$_2$. This report reviews the theory
of each extension, describes how MRST implements it, and documents a
complete reproduction of the published examples. All runs used GNU
Octave~8.4 instead of MATLAB, on the MRST development sources and the
example decks published with the paper. The reproduction therefore also
serves as a portability and reproducibility study
(Section~\ref{sec:practice}).

Section~\ref{sec:model} states the governing equations and the
discretization. Section~\ref{sec:trapping} presents the spill-point
analysis of structural trapping. Sections~\ref{sec:pvt}--\ref{sec:diff}
cover the three extensions and their verification cases.
Section~\ref{sec:practice} reports computational aspects and the
portability findings. Section~\ref{sec:conclusions} concludes.
Appendix~\ref{app:wellindex} audits the well-index model of
Section~\ref{sec:diff} against a grid-aware alternative for unstructured
meshes.

\section{Governing equations and discretization}
\label{sec:model}

\subsection{Black-oil formulation for CO$_2$--brine systems}

The black-oil model describes a multiphase system with pseudo-components
whose mutual solubility depends on pressure only. For GCS in saline
aquifers, the oleic phase is assigned brine properties, so that the
gaseous (CO$_2$-rich) pseudo-component can dissolve in it
\cite{hassanzadeh2008,salo2024}. With subscripts $o$ (aqueous/oleic) and
$g$ (gaseous), mass conservation of the two pseudo-components reads
\begin{align}
\frac{\partial}{\partial t}\Big[\phi\big(b_o S_o + R_v b_g S_g\big)\Big]
 + \nabla\cdot\big(b_o \mathbf{u}_o + R_v b_g \mathbf{u}_g\big) &= q_o,
\label{eq:bo-oleic}\\
\frac{\partial}{\partial t}\Big[\phi\big(b_g S_g + R_s b_o S_o\big)\Big]
 + \nabla\cdot\big(b_g \mathbf{u}_g + R_s b_o \mathbf{u}_o\big) &= q_g,
\label{eq:bo-gas}
\end{align}
where $\phi$ is porosity, $S_\alpha$ saturation,
$b_\alpha = 1/B_\alpha$ the reciprocal formation volume factor, $R_s$ the
solution gas--oil ratio (dissolved CO$_2$ in brine), $R_v$ the vaporized
oil--gas ratio (H$_2$O in the CO$_2$-rich phase), and $q_\alpha$ source
terms. The phase fluxes follow the multiphase extension of Darcy's law,
\begin{equation}
\mathbf{u}_\alpha = -\frac{\mathbf{k}\, k_{r\alpha}(S_\alpha)}
{\mu_\alpha}\big(\nabla p_\alpha - \rho_\alpha \mathbf{g}\big),
\label{eq:darcy}
\end{equation}
with intrinsic permeability $\mathbf{k}$, relative permeability
$k_{r\alpha}$, viscosity $\mu_\alpha$, and phase pressures related by the
capillary pressure $p_g - p_o = P_c(S_g)$. The saturations close the
system through $S_o + S_g = 1$. Reservoir-condition phase densities
follow from the surface densities $\rho_\alpha^s$:
\begin{equation}
\rho_o = b_o\big(\rho_o^s + R_s\,\rho_g^s\big), \qquad
\rho_g = b_g\big(\rho_g^s + R_v\,\rho_o^s\big).
\label{eq:densities}
\end{equation}

\subsection{Discretization and solvers in MRST}

MRST discretizes \eqref{eq:bo-oleic}--\eqref{eq:darcy} with a
cell-centered finite-volume method: two-point flux approximation (TPFA)
for the pressure gradient, single-point upstream mobility weighting, and
backward Euler in time \cite{lie2019}. For a face $ij$ between cells $i$
and $j$, the discrete phase flux is
\begin{equation}
v_{\alpha,ij} = \lambda_{\alpha}^{\mathrm{up}}\, T_{ij}
\big(p_{\alpha,i} - p_{\alpha,j} - \rho_{\alpha,ij}\, g\, \Delta z_{ij}\big),
\qquad
T_{ij} = \Big(T_{i,ij}^{-1} + T_{j,ij}^{-1}\Big)^{-1},
\label{eq:tpfa}
\end{equation}
where $T_{i,ij}$ are half-face transmissibilities built from the cell
permeability and geometry, and
$\lambda_\alpha^{\mathrm{up}} = (k_{r\alpha}/\mu_\alpha)^{\mathrm{up}}$
is evaluated in the upstream cell. The fully implicit residual system is
linearized by Newton's method. MRST assembles the Jacobians by automatic
differentiation and organizes secondary variables (densities, mobilities,
relative permeabilities, component fluxes) as \emph{state functions}
attached to the model \cite{lie-moyner2021}. This design is what allows
the hysteresis and diffusion extensions of
Sections~\ref{sec:hyst}--\ref{sec:diff} to replace one state function
each, without changes to the assembly loop.

In this work all systems were solved with a sparse direct solver
(\texttt{BackslashSolverAD}), a Newton line search, and an
iteration-targeting time-step controller. The Octave environment provides
no compiled (MEX) acceleration; Section~\ref{sec:practice} quantifies the
consequences.

\section{Structural trapping: spill-point analysis}
\label{sec:trapping}

Before dynamic simulation, geometric analysis of the caprock surface
already quantifies the structural trapping potential. The
\texttt{co2lab} module builds a top-surface grid $\Gamma$ from the 3D
grid and analyzes its depth function $z(\mathbf{x})$
\cite{nilsen2015}. A \emph{trap} is the interior of a closed depth
contour around a local minimum of $z$ (a dome apex); its \emph{spill
point} is the saddle where CO$_2$ overflows to the next trap. The traps
and the connecting \emph{spill paths} form a directed tree that predicts
the long-term migration of buoyant CO$_2$ without solving any flow
equation.

We exercised this analysis on the introductory example of
\texttt{co2lab-spillpoint}: a synthetic $10\times5$~km, 50~m thick
sandbox at 1~km depth, with a dipping, sinusoidally perturbed top
surface (100$\times$100 top-surface cells, porosity 0.25). The
edge-based algorithm of \texttt{trapAnalysis} identified 13 traps with a
total capacity of \SI{5.099e6}{m^3}, which is 40.8\% of the total pore
volume (\SI{1.25e7}{m^3}); the five largest traps hold 55.2\% of that
capacity. Figure~\ref{fig:traps} shows the depth map with the trap cells
and the trap--spill-path network.

\begin{figure}[tb]
\centering
\includegraphics[width=0.495\textwidth]{fig1_depth_and_traps}
\includegraphics[width=0.495\textwidth]{fig2_traps_and_spillpaths}
\caption{Spill-point analysis of the synthetic sandbox aquifer.
Left: top-surface depth (m); white dots mark cells inside structural
traps. Right: the 13 numbered traps (red), the spill paths connecting
them (green), and the remaining cells (blue). CO$_2$ that leaks from a
trap migrates along the spill path to the next trap up-dip.}
\label{fig:traps}
\end{figure}

\section{PVT data generation for CO$_2$--brine black-oil models}
\label{sec:pvt}

\subsection{Thermodynamic model}

Black-oil simulators consume tabulated PVT data. For CO$_2$--brine
systems, the tables must encode the mutual solubility of CO$_2$ and
H$_2$O and the resulting density and viscosity changes. The
implementation of \cite{salo2024} follows Hassanzadeh et
al.~\cite{hassanzadeh2008}: phase compositions come from the
fugacity--activity equilibrium model of Spycher et
al.~\cite{spycher2003,spycher2005} combined with the activity
coefficients of Duan and Sun~\cite{duan2003}, which admits NaCl, KCl,
CaCl$_2$, and seawater brines; the model is valid to approximately
\SI{100}{\celsius} and \SI{600}{bar}. Gas-phase molar volumes follow the
Redlich--Kwong equation of state \cite{redlich1949}. Brine density
follows Rowe and Chou \cite{rowe1970}, extended in $p$--$T$ range by
Batzle and Wang \cite{batzle1992}, with the density increase due to
dissolved CO$_2$ from Garcia's apparent molar volume correlation
\cite{garcia2001}. Viscosities use Islam and Carlson \cite{islam2012}
for the aqueous phase and Fenghour et al.~\cite{fenghour1998} for
CO$_2$, with Davidson's mixing rule \cite{davidson1993} when the gas
phase contains vaporized H$_2$O.

The black-oil quantities follow from the equilibrium compositions. The
solution gas--oil ratio $R_s$ is the volume of dissolved CO$_2$,
measured at surface conditions, per surface volume of brine, and $B_o$
is the ratio of the reservoir-condition aqueous volume to its surface
volume. If vaporized water is included, the gas-phase analogues read
\cite{salo2024}
\begin{equation}
R_v \;=\; \frac{V^s_{v\mathrm{H_2O}}}{V^s_{\mathrm{CO_2}}}
      \;=\; \frac{y_{\mathrm{H_2O}}\,\bar\rho^{\,s}_{\mathrm{CO_2}}}
                 {\bar\rho^{\,s}_{\mathrm{H_2O}}\, y_{\mathrm{CO_2}}},
\qquad
B_g \;=\; \frac{V^r_{\mathrm{g}}}{V^s_{\mathrm{CO_2}}}
      \;=\; \frac{\bar\rho^{\,s}_{\mathrm{CO_2}}}
                 {\bar\rho^{\,r}_{\mathrm{g}}\,(1 - v_{\mathrm{H_2O}})},
\label{eq:rvbg}
\end{equation}
where $y$ and $v$ are mole and mass fractions in the gaseous phase and
$\bar\rho$ denotes molar density. The function
\texttt{pvtBrineWithCO2BlackOil} evaluates this chain over a requested
$p$--$T$--salinity range and writes ECLIPSE-format \texttt{PVTO} (live
oil, i.e., brine with dissolved CO$_2$) and \texttt{PVDG} (dry gas)
tables.

\subsection{Verification}

The published deck \texttt{example\_co2brine\_3reg\_unsatVals.DATA}
embeds tables generated at $T=\SI{80}{\celsius}$, salinity $10^5$~ppm,
and $p \in [1,600]$~bar. Reading this deck with
\texttt{initDeckADIFluid} produces the MRST fluid object; evaluating its
$b_g$, $\mu_g$, $R_s^{\mathrm{sat}}$, $b_o$, $\mu_o$ functions and
reconstructing densities through \eqref{eq:densities} must return the
tabulated values. Figure~\ref{fig:pvt} shows the comparison; the maximum
relative difference over all four quantities is $2.2\times 10^{-16}$,
i.e., machine precision, which verifies the deck round-trip
(table $\rightarrow$ deck $\rightarrow$ interpolated fluid object). The
curves also show the physics of interest for GCS: the supercritical
density transition of CO$_2$ (\num{223}, \num{734}, and
\SI{933}{kg/m^3} at 100, 300, and 600~bar) and the density increase of
brine upon CO$_2$ dissolution, which drives the convective mixing of
Section~\ref{sec:diff}.

\begin{figure}[tb]
\centering
\includegraphics[width=\textwidth]{fig_pvt_props}
\caption{Verification of the PVT deck round-trip at
$T=\SI{80}{\celsius}$, $S=10^5$~ppm. Thick lines: MRST fluid object;
thin lines: tabulated output of \texttt{pvtBrineWithCO2BlackOil}. From
left to right: CO$_2$ density, CO$_2$ viscosity (log scale), brine
density without and with dissolved CO$_2$, brine viscosity. The curves
coincide to machine precision; compare with Fig.~2 of \cite{salo2024}.}
\label{fig:pvt}
\end{figure}

\section{Relative permeability hysteresis}
\label{sec:hyst}

\subsection{Killough's model with Land trapping}

Relative permeability depends on the saturation path, not only on the
current saturation: pore-scale contact-angle hysteresis and nonwetting
phase disconnection make the imbibition curve differ from the drainage
curve \cite{juanes2006}. In GCS, CO$_2$ is the nonwetting phase in most
aquifers; during post-injection migration, brine re-imbibes at the
trailing edge of the plume, disconnects CO$_2$ ganglia, and immobilizes
them (residual trapping). A drainage-only model misses this mechanism
and overpredicts the mobile gas cap.

Killough's model \cite{killough1976} computes the scanning curve between
the primary drainage curve $k^d_{rg}$ and the bounding imbibition curve
$k^{ib}_{rg}$. While the current saturation $S_g$ lies below the
historical maximum $S_{gi}$, the scanning relative permeability is
\begin{equation}
k^{i}_{rg}(S_g) = k^{ib}_{rg}(S_g^*)\,
\frac{k^{d}_{rg}(S_{gi})}{k^{d}_{rg}(S_{g,\max})},
\qquad
S_g^* = S_{gt,\max} + \frac{(S_g - S_{gt})(S_{g,\max} -
S_{gt,\max})}{S_{gi} - S_{gt}},
\label{eq:killough}
\end{equation}
where $S_{g,\max}$ is the maximum attainable gas saturation and
$S_{gt,\max}$ the corresponding maximum trapped saturation. The trapped
gas saturation for the current cycle follows Land's model
\cite{land1968},
\begin{equation}
S_{gt} = S_{g,\min} + \frac{S_{gi} - S_{g,\min}}
{1 + C\,(S_{gi} - S_{g,\min})},
\qquad
C = \frac{1}{S_{gt,\max} - S_{g,\min}} - \frac{1}{S_{g,\max} -
S_{g,\min}},
\label{eq:land}
\end{equation}
with $S_{g,\min}=0$ on primary drainage. For numerical robustness the
implementation replaces the denominator $1 + C(S_{gi}-S_{g,\min})$ by
$A + C(S_{gi}-S_{g,\min})$ with $A = 1 + a(S_{g,\max}-S_{gi})$ and
default $a=0.1$, as ECLIPSE does \cite{salo2024}. MRST reads the
hysteresis options from the \texttt{EHYSTR} keyword and the imbibition
saturation tables from the deck, and evaluates \eqref{eq:killough} inside
a dedicated state function, \texttt{HystereticRelativePermeability},
that replaces the default relative permeability of the validated model.
Figure~\ref{fig:kr} shows the three curves for the PUNQ-S3 rock (Berea
sandstone data of Oak, used by \cite{juanes2006}), computed from the
published deck.

\begin{figure}[tb]
\centering
\includegraphics[width=0.55\textwidth]{fig_kr_hysteresis}
\caption{Hysteretic gas relative permeability of the PUNQ-S3 deck:
primary drainage curve, bounding imbibition curve, and the Killough
scanning curve for a flow reversal at $S_{gi}=0.4$. The scanning curve
reaches zero at the Land trapped saturation
$S_{gt}(S_{gi})\approx 0.29$. Compare with Fig.~3 of \cite{salo2024}.}
\label{fig:kr}
\end{figure}

\subsection{Validation on the PUNQ-S3 model}

Juanes et al.~\cite{juanes2006} adapted the PUNQ-S3 geologic model into a
GCS benchmark: a $19\times28\times5$ corner-point grid (1{,}761 active
cells) of an anticline aquifer with heterogeneous permeability
(0.5--1000~mD), eight injectors at \SI{18}{rm^3/day} each for 10 years,
500 years of total simulation, immiscible water--gas physics, and pore
volumes multiplied by $10^3$ on the boundary ring to mimic an open
aquifer. We reproduced both published cases from the same deck
(\texttt{CASE2.DATA}): Case~1 deactivates hysteresis; Case~2 activates
Killough hysteresis in the gas phase.

Figure~\ref{fig:punq-maps} shows the CO$_2$ saturation in all five
layers at $t=500$~y, and Fig.~\ref{fig:punq-evol} the plume metrics in
time. The two cases coincide until injection stops at $t=10$~y, because
hysteresis only acts on saturation reversals. Afterwards they diverge:
without hysteresis the plume migrates up-dip, drains the invaded region,
and collects into a compact, high-saturation gas cap (49 cells with
$S_g>0.05$, mean plume saturation 0.505 at 500~y); with hysteresis the
trailing edge traps CO$_2$ residually along the migration path, the
invaded area freezes at its maximum extent (138 cells), and the mean
plume saturation drops to 0.192. Table~\ref{tab:punq} summarizes the
metrics. This is precisely the behavior reported in Fig.~4 of
\cite{salo2024} and in \cite{juanes2006}, where the same contrast was
validated against ECLIPSE~100.

\begin{figure}[tb]
\centering
\includegraphics[width=\textwidth]{fig_punq_sg500y}
\caption{PUNQ-S3: CO$_2$ saturation per layer (left to right: layers
1--5, top of the formation to its base) at $t=500$~y. Top row: no
hysteresis; bottom row: Killough hysteresis. Gray marks inactive cells.
Hysteresis suppresses the compact gas cap and leaves a broad,
low-saturation, residually trapped plume. Compare with Fig.~4 of
\cite{salo2024}.}
\label{fig:punq-maps}
\end{figure}

\begin{figure}[tb]
\centering
\includegraphics[width=\textwidth]{fig_punq_evolution}
\caption{PUNQ-S3 plume evolution. Left: number of cells with
$S_g>0.05$; right: maximum gas saturation. The curves coincide during
injection ($t\le 10$~y) and diverge after flow reversal, when the
hysteretic scanning curves activate.}
\label{fig:punq-evol}
\end{figure}

\begin{table}[tb]
\centering
\caption{PUNQ-S3 plume metrics at $t=500$~y.}
\label{tab:punq}
\begin{tabular}{lcc}
\toprule
Metric & No hysteresis & Hysteresis \\
\midrule
Maximum $S_g$                & 0.690 & 0.641 \\
Cells with $S_g>0.05$        & 49    & 138   \\
Cells with $S_g>0.3$         & 36    & 23    \\
Mean $S_g$ where $S_g>0.05$  & 0.505 & 0.192 \\
\bottomrule
\end{tabular}
\end{table}

\section{Molecular diffusion of dissolved CO$_2$}
\label{sec:diff}

\subsection{Model and numerical diffusion}

When CO$_2$ dissolves into brine, the mixture is denser than the
resident brine and sinks in Rayleigh--Taylor fingers, which accelerates
dissolution trapping. Concentration gradients also drive a diffusive
mass flux. The total macroscale flux of component $\gamma$ in phase
$\alpha$ decomposes as \cite{bear1972,salo2024}
\begin{equation}
\mathbf{J}^\gamma_\alpha =
\underbrace{\rho_\alpha\,\chi^\gamma_\alpha\,\mathbf{u}_\alpha}_{\text{advective}}
\;-\;
\underbrace{\phi\,\rho_\alpha S_\alpha\, D^\gamma_\alpha\,
\nabla\chi^\gamma_\alpha}_{\text{diffusive (Fickian)}},
\label{eq:diffflux}
\end{equation}
where $\chi^\gamma_\alpha$ is the mass fraction and $D^\gamma_\alpha$ a
scalar pseudo-diffusivity that lumps molecular diffusivity and
tortuosity; for CO$_2$ in brine the molecular value is
$\mathcal{D}\sim 2$--$5\times10^{-9}$~m$^2$/s. MRST implements
\eqref{eq:diffflux} by replacing the \texttt{ComponentTotalFlux} state
function of \texttt{GenericBlackOilModel} with
\texttt{CO2TotalFluxWithDiffusion}, which adds the second term with a
TPFA-style face diffusivity.

Whether the diffusive term matters numerically depends on the grid. A
first-order upstream scheme carries a numerical diffusion
$D_{\mathrm{num}}\sim u\,h$, with $u$ the Darcy velocity and $h$ the
cell size. For a laboratory-scale problem with $k=100$~mD,
$\mu=0.8$~cP, and a pressure gradient of 10~mbar/m,
$u\approx1.2\times10^{-7}$~m/s, so $h=1$~cm gives
$D_{\mathrm{num}}\sim10^{-9}$~m$^2$/s---larger than the physical
diffusivity---while $h=2.5$~mm brings the two within reach of each
other. Physical diffusion therefore only shapes the solution when the
grid is fine enough, and conversely, coarse-grid simulations already
contain more (numerical) diffusion than the physics requires
\cite{salo2024}.

\subsection{A FluidFlower-inspired convective mixing problem}

The test case is a quasi-2D stratigraphic section
($1\times0.66$~m, one 1-cm-thick cell layer) inspired by the FluidFlower
benchmark rig \cite{flemisch2024}, repositioned at 1~km depth so that
CO$_2$ is supercritical. The PEBI mesh conforms to a fault structure
(fault permeability 10~mD inside a 1~D reservoir); sealing units are
removed from the flow domain, and lateral boundary cells receive a
$10^5$ pore-volume multiplier to emulate an open aquifer. CO$_2$ is
injected in the lower-right reservoir at \SI{8}{mL/min} (surface
conditions) for one day, with rate ramps, and the model then runs to
$t=30$~days with dissolution active (\texttt{disgas}).

We ran the case on two meshes: coarse ($h\approx1$~cm, 8{,}206 cells)
and fine ($h\approx2.5$~mm, 105{,}279 cells, the resolution of Fig.~6 in
\cite{salo2024}), each with $D=0$ and $D=2\times10^{-11}$~m$^2$/s.
Figure~\ref{fig:diff-coarse} shows the dissolved CO$_2$ field
($R_s$) at $t=30$~d on the coarse mesh: a CO$_2$-saturated gas
accumulation under the sealing horizon, and sinking CO$_2$-rich brine
below it. At this resolution the two diffusivities give nearly identical
fields (mean $R_s$ 4.06 vs.\ 4.13~Sm$^3$/Sm$^3$; plume cell count 1{,}027
vs.\ 1{,}105), which confirms the scaling argument: numerical diffusion
at $h=1$~cm dominates the physical value.
%
Figure~\ref{fig:diff-fine} repeats the comparison on the fine mesh,
where the convective fingers are resolved and the physical diffusivity
becomes visible, as in Fig.~6 of \cite{salo2024}: without diffusion the
fingers below the gas accumulation are thin and closely spaced; with
$D=2\times10^{-11}$~m$^2$/s they are wider, smoother, and farther
apart. Diffusion also enhances dissolution trapping: the mean dissolved
ratio rises from 2.97 to 3.54~Sm$^3$/Sm$^3$, and the invaded region
grows from 9{,}574 to 10{,}389 cells.
\begin{figure}[tb]
\centering
\includegraphics[width=\textwidth]{fig_diffusion_fine30d}
\caption{Fine mesh ($h\approx2.5$~mm, 105{,}279 cells): dissolved CO$_2$
($R_s$) at $t=30$~d without (top) and with (bottom) molecular diffusion
($D=2\times10^{-11}$~m$^2$/s). Finger width and spacing increase with
the diffusivity. Compare with Fig.~6(a,b) of \cite{salo2024}.}
\label{fig:diff-fine}
\end{figure}

\begin{figure}[tb]
\centering
\includegraphics[width=\textwidth]{fig_diffusion_rs30d}
\caption{Coarse mesh ($h\approx1$~cm, 8{,}206 cells): dissolved CO$_2$
($R_s$, Sm$^3$/Sm$^3$) at $t=30$~d, without (top) and with (bottom)
molecular diffusion. Dark: CO$_2$-saturated brine under the gas
accumulation; the halo below is the sinking CO$_2$-rich brine. At this
resolution numerical diffusion masks the physical diffusivity.}
\label{fig:diff-coarse}
\end{figure}

\section{Computational practice and portability}
\label{sec:practice}

\subsection{Runtimes}

All simulations ran fully implicitly with a sparse direct linear solver,
Newton line search, and iteration-targeting time steps, in GNU
Octave~8.4 on a Linux workstation, single-threaded, without any MEX
acceleration. Table~\ref{tab:runtimes} lists the runtimes. For
reference, \cite{salo2024} reports \SI{6857}{s} for the 90{,}000-cell
Johansen field case in MATLAB with a direct solver, reduced to
\SI{417}{s} with a compiled AD backend and an algebraic multigrid CPR
solver; that case was out of scope here for exactly this reason.

\begin{table}[tb]
\centering
\caption{Reproduction runtimes (GNU Octave 8.4, single thread,
direct linear solver).}
\label{tab:runtimes}
\begin{tabular}{lrrr}
\toprule
Case & Cells & Report steps & Runtime [s] \\
\midrule
PUNQ-S3, no hysteresis      & 1{,}761   & 100 & 33.5 \\
PUNQ-S3, hysteresis         & 1{,}761   & 100 & 51.4 \\
Diffusion, coarse, $D{=}0$    & 8{,}206   & 69  & 167.9 \\
Diffusion, coarse, $D{=}2\times10^{-11}$ & 8{,}206 & 69 & 166.9 \\
Diffusion, fine, $D{=}0$      & 105{,}279 & 69 & 3531.7 \\
Diffusion, fine, $D{=}2\times10^{-11}$   & 105{,}279 & 69 & 2751.0 \\
\bottomrule
\end{tabular}
\end{table}

\subsection{MATLAB-to-Octave portability}

The reproduction required a set of well-defined substitutions, all
outside the numerical kernels: the \texttt{table} class (used only by
the PVT table writer; the shipped tables were used instead),
\texttt{deval} (replaced by explicit midpoint integration of the
hydrostatic ODE), \texttt{pdist2} (replaced by a direct norm),
\texttt{griddedInterpolant} (MRST ships an Octave replacement of
\texttt{fastInterpTable} that must shadow the MATLAB version), and
\texttt{delaunayTriangulation} (a minimal shim exposing only the
\texttt{edges} method suffices for the PEBI grid generator).

One finding is of general interest for simulation-software
reproducibility. The PEBI meshes of Section~\ref{sec:diff} initially
came out with 38\% of cells having negative volume, without any error.
The cause was a malformed block comment in
\texttt{sortEdges.m} of the \texttt{upr} module: the \texttt{\%\}}
closer was glued to the end of the copyright line. MATLAB tolerates
this; Octave parses the block as unterminated, silently treats the whole
function body as a comment, and the function degenerates to an identity
map. The unsorted cell topology then produced inverted prisms during
extrusion---three calls downstream of the defect. A second file
(\texttt{mrstGridWithFullMappings.m}) lost its entire function
definition to the same pattern. The defect exists only in MRST's subtree
copy (the upstream UPR repository is well-formed), which points to a
mass copyright-update script as the origin. The fix was submitted
upstream.\footnote{\url{https://github.com/SINTEF-AppliedCompSci/MRST/pull/37}}
The episode illustrates two practical lessons: (i) cross-interpreter
parsing differences can corrupt \emph{results} rather than raise errors,
so geometric invariants (here: positive cell volumes and face areas) are
worth asserting after every grid-construction step; and (ii) automated
rewriting of source files, however trivial, deserves a syntax check in
all supported interpreters.

\subsection{Companion Jupyter notebooks}

The reproduction is also released as four executable Jupyter notebooks
(\texttt{reproductions/salo2024-gcs/notebooks/}), one per case
(structural trapping, PVT, hysteresis, diffusion), running MRST through
an Octave Jupyter kernel with markdown narrative paraphrased from this
report. The expensive cases (hysteresis, diffusion) default to loading
the cached \texttt{results/*.mat} states rather than re-simulating, with
a boolean flag to opt into a live re-run; both paths were verified to
reproduce the numbers in this report exactly. Building them surfaced one
further portability pitfall, specific to the interactive kernel rather
than to headless scripts: this Octave build has no \texttt{qt} graphics
toolkit, and \texttt{fltk} is needed for the 3-D patch plots (\texttt{plotGrid}
on a 3-D grid) that gnuplot cannot draw; but forcing \texttt{fltk} as the
session-wide default silently renders every \texttt{imagesc} plot --
the type used throughout this report's own figures -- as a blank raster,
with no error. The fix is to keep gnuplot as the default and switch to
\texttt{fltk} only in the specific cell that draws a 3-D patch plot.

\section{Conclusions}
\label{sec:conclusions}

We reproduced, in GNU Octave, the three \texttt{ad-blackoil} extensions
for geologic CO$_2$ storage published by Sal\'o-Salgado et
al.~\cite{salo2024}, plus a spill-point structural trapping analysis.
The PVT deck round-trip is exact to machine precision. The PUNQ-S3
validation reproduces the published effect of relative permeability
hysteresis: residual trapping freezes the plume at its maximum extent
and suppresses the mobile gas cap, with divergence starting exactly at
flow reversal. The convective mixing example reproduces both the physics
(dissolution-driven density fingering) and the numerical analysis: the
physical diffusivity only matters once the grid-level numerical
diffusion drops below it. The full reproduction, including a 105k-cell
fine-mesh run, is feasible on open-source tooling in minutes to tens of
minutes per case; the exercise also uncovered and fixed a silent
grid-corruption defect, which underlines the value of independent
reproduction across interpreters. Appendix~\ref{app:wellindex} shows that
MRST's well-index model on this PEBI mesh is a Cartesian bounding-box
approximation of Peaceman's formula, not a grid-aware one: the gap against
a Voronoi-specific alternative grows with cell irregularity, though it
stays within this problem's large injectivity margin.

\appendix

\section{Peaceman vs.\ Palagi--Aziz well index on a PEBI mesh}
\label{app:wellindex}

Section~\ref{sec:model} characterized MRST's well-reservoir coupling as
Peaceman-type~\cite{peaceman1978}. This appendix examines that
characterization on the diffusion case's own mesh
(Section~\ref{sec:diff}), whose injector cell is a PEBI polygon, not a
Cartesian block, and quantifies the effect of a grid-aware alternative.

\subsection{The coupling and well index used by MRST}

For a single perforated cell, the rate/pressure relation and well index
implemented in \texttt{core/utils/computeWellIndex.m} are
\begin{equation}
q = -WI\cdot\mathrm{mob}\cdot\big(p_{\mathrm{cell}} - p_{bh} - \Delta
p_{\mathrm{hydrostatic}}\big), \qquad
WI = \frac{2\pi\sqrt{k_x k_y}\,h}{\ln(r_0/r_w) + \mathrm{skin}},
\label{eq:wi-peaceman}
\end{equation}
with the anisotropic equivalent radius of Peaceman~\cite{peaceman1983}
\begin{equation}
r_0 = \frac{0.28\sqrt{\sqrt{k_y/k_x}\,\Delta x^2 +
\sqrt{k_x/k_y}\,\Delta y^2}}{(k_y/k_x)^{1/4}+(k_x/k_y)^{1/4}}.
\label{eq:r0-peaceman}
\end{equation}
Equation~\eqref{eq:r0-peaceman} assumes the well cell \emph{is} a
rectangle of known $\Delta x,\Delta y$ -- true by construction on a
Cartesian or corner-point grid, but undefined for an irregular PEBI
polygon.

\subsection{The bounding-box approximation on a PEBI grid}

For any grid with a node list, MRST's \texttt{cellDims.m} supplies
$\Delta x,\Delta y$ to \texttt{computeWellIndex} as the axis-aligned
bounding box of the well cell's own nodes, whatever its true shape: it
takes the minimum and maximum node coordinates and sets $\Delta x = M_x -
m_x$, $\Delta y = M_y - m_y$. This never inspects the well cell's actual
neighbors, connection transmissibilities, or number of sides -- two very
differently shaped cells with the same bounding box receive the same well
index.

\subsection{Palagi and Aziz's well index for Voronoi grids}

Palagi and Aziz~\cite{palagi1994a,palagi1994b} derive an equivalent
radius directly from the well cell's actual connections, so it is defined
for any polygon. Assuming radial flow around the well-cell gridpoint out
to each neighbor $j$,
\begin{equation}
p_j - p_w = \frac{qB\mu\,\ln(L_{ij}/r_w)}{\theta kh},
\label{eq:pa-radial}
\end{equation}
and requiring this to be consistent with the discretized single-phase
flow equation for the well cell,
\begin{equation}
q = \sum_j \frac{T_{ij}}{B\mu}(p_j - p_o)
  = \sum_j \frac{T_{ij}}{B\mu}(p_j - p_w) - \sum_j
    \frac{T_{ij}}{B\mu}(p_o - p_w),
\label{eq:pa-discrete}
\end{equation}
gives the equivalent radius as a transmissibility-weighted average over
the real connections,
\begin{equation}
\ln r_{eq} = \frac{\sum_j T_{ij}\ln L_{ij} - \theta kh}{\sum_j T_{ij}},
\qquad
WI = \frac{\theta kh}{\ln(r_{eq}/r_w) + \mathrm{skin}},
\label{eq:pa-req}
\end{equation}
where $T_{ij}$ is the two-point transmissibility of connection $ij$,
$L_{ij}$ the distance between cell centroids, and $\theta$ the angle open
to flow ($\theta=2\pi$ for a fully interior well cell, matching the
implicit $2\pi$ in Eq.~\eqref{eq:wi-peaceman}, which keeps the two models
directly comparable). Equation~\eqref{eq:pa-req} reduces to Peaceman's
formula for a square Cartesian cell but stays correct for an irregular
polygon with any number of neighbors. We implemented
Eq.~\eqref{eq:pa-req} as a standalone function
(\texttt{scripts/computeWellIndexPalagiAziz.m}) that reuses MRST's own
\texttt{cellDims} and \texttt{computeTrans}, matching $kh$, skin, and
$r_w$ exactly with Eq.~\eqref{eq:wi-peaceman} so the two well indices
differ only in how $r_{eq}$ is derived; \texttt{addWell} accepts a
precomputed \texttt{WI} directly, so no MRST internals needed modifying.

\subsection{Comparison}

Both models were evaluated at $r_w=1$~mm, skin $=0$, on the coarse mesh
of Section~\ref{sec:diff}, at two locations: the paper's actual injection
cell (in the reservoir matrix, nearly a cube,
$\Delta x\approx\Delta y\approx\Delta z\approx1$~cm), and the most
elongated fault-strip cell reachable near that same location (bounding-box
aspect ratio $R_y\approx2.4$, found by ranking all fault cells with
\texttt{scripts/find\_fault\_cell.m}) -- Palagi and Aziz~\cite{palagi1994b}
flag $R_y>2$ as exactly the regime where the Cartesian formula stops being
reliable. Table~\ref{tab:wi} shows the well-index gap growing with cell
irregularity, as expected.

\begin{table}[tb]
\centering
\caption{Well index at $r_w=1$~mm, skin $=0$: Peaceman (MRST default)
vs.\ Palagi--Aziz (Eq.~\ref{eq:pa-req}).}
\label{tab:wi}
\begin{tabular}{lccc}
\toprule
Location & $WI$ Peaceman & $WI$ Palagi--Aziz & Difference \\
\midrule
Reservoir well (original)  & $9.034\times10^{-14}$ & $8.157\times10^{-14}$ & $-9.71\%$ \\
Elongated fault-strip cell & $5.555\times10^{-16}$ & $4.583\times10^{-16}$ & $-17.50\%$ \\
\bottomrule
\end{tabular}
\end{table}

Because this well is rate-controlled, $WI$ only sets the drawdown
$p_{\mathrm{cell}}-p_{bh}$ needed to deliver the fixed 8~mL/min --
it never changes how much CO$_2$ enters the reservoir, so the
saturation and dissolved-CO$_2$ fields should be governed by rate, not by
$WI$, and only the bottom-hole pressure should respond. We reran the
coarse-mesh, $D=0$ case with each well index at both locations
(Table~\ref{tab:wi-effect}): the field-wide difference stays at the level
of solver noise at both locations, but the bottom-hole-pressure gap grows
about $230\times$ from the reservoir well to the fault-strip cell,
tracking the larger well-index gap and the fault's lower permeability
(10~mD vs.\ 1000~mD, i.e.\ less injectivity headroom).

\begin{table}[tb]
\centering
\caption{Effect of the well-index model on the coarse-mesh, $D=0$
simulation (max.\ absolute difference over the full 30-day schedule).}
\label{tab:wi-effect}
\begin{tabular}{lccc}
\toprule
Location & Max $|\Delta p_{bh}|$ [Pa] & Max $|\Delta r_s|$ & Max $|\Delta S_g|$ \\
\midrule
Reservoir well              & 0.35 & $1.6\times10^{-6}$ & $3.8\times10^{-9}$ \\
Elongated fault-strip cell  & 80.0 & $6.5\times10^{-6}$ & $1.5\times10^{-7}$ \\
\bottomrule
\end{tabular}
\end{table}

\subsection{Discussion}

MRST's well model is Peaceman-type, but the equivalent radius it computes
for a PEBI cell is a bounding-box stand-in for the Cartesian formula, not
a generalization of it. That gap is real and grows with cell irregularity,
here about twice as large at an elongated fault-strip cell as at a
well-conditioned reservoir cell. Whether it matters for a given simulation
is a separate question: for this rate-controlled injector it changes only
the computed bottom-hole pressure, never the injected mass, and even the
largest pressure gap found (80~Pa against a $\sim$9.9~MPa reservoir
pressure) is negligible at this problem's flow rate -- the reservoir has
far more injectivity headroom than this well ever calls on. A
pressure-controlled well, a higher rate, or a lower-permeability target
would be expected to show a larger effect, since drawdown itself would no
longer be negligible.

\section*{Acknowledgments}

This report and the reproduction runs were prepared with the assistance
of Claude Code (Anthropic). MRST is developed by the Applied
Computational Science group at SINTEF Digital. The example decks and
code are from L.~Sal\'o-Salgado's \texttt{mrst-adblackoil-gcs}
repository and MRST's \texttt{co2lab-mit} module.

\begin{thebibliography}{00}

\bibitem{benson2005} S.~Benson, P.~Cook, et al., Underground geological
storage, in: IPCC Special Report on Carbon Dioxide Capture and Storage,
Cambridge University Press, 2005, pp. 196--276.

\bibitem{lie2019} K.-A. Lie, An Introduction to Reservoir Simulation
Using MATLAB/GNU Octave: User Guide for the MATLAB Reservoir Simulation
Toolbox (MRST), Cambridge University Press, 2019.

\bibitem{lie-moyner2021} K.-A. Lie, O.~M{\o}yner (Eds.), Advanced
Modelling with the MATLAB Reservoir Simulation Toolbox, Cambridge
University Press, 2021.

\bibitem{salo2024} L.~Sal\'o-Salgado, O.~M{\o}yner, K.-A. Lie,
R.~Juanes, Three-dimensional simulation of geologic carbon dioxide
sequestration using MRST, Advances in Geo-Energy Research 14 (2024)
34--48. \url{https://doi.org/10.46690/ager.2024.10.06}

\bibitem{hassanzadeh2008} H.~Hassanzadeh, M.~Pooladi-Darvish,
A.~M. Elsharkawy, D.~W. Keith, Y.~Leonenko, Predicting PVT data for
CO$_2$--brine mixtures for black-oil simulation of CO$_2$ geological
storage, International Journal of Greenhouse Gas Control 2 (2008)
65--77.

\bibitem{nilsen2015} H.~M. Nilsen, K.-A. Lie, O.~Andersen, Analysis of
CO$_2$ trapping capacities and long-term migration for geological
formations in the Norwegian North Sea using MRST-co2lab, Computers \&
Geosciences 79 (2015) 15--26.

\bibitem{spycher2003} N.~Spycher, K.~Pruess, J.~Ennis-King,
CO$_2$--H$_2$O mixtures in the geological sequestration of CO$_2$. I.
Assessment and calculation of mutual solubilities from 12 to 100$^\circ$C
and up to 600 bar, Geochimica et Cosmochimica Acta 67 (2003) 3015--3031.

\bibitem{spycher2005} N.~Spycher, K.~Pruess, CO$_2$--H$_2$O mixtures in
the geological sequestration of CO$_2$. II. Partitioning in chloride
brines at 12--100$^\circ$C and up to 600 bar, Geochimica et Cosmochimica
Acta 69 (2005) 3309--3320.

\bibitem{duan2003} Z.~Duan, R.~Sun, An improved model calculating CO$_2$
solubility in pure water and aqueous NaCl solutions from 273 to 533 K
and from 0 to 2000 bar, Chemical Geology 193 (2003) 257--271.

\bibitem{redlich1949} O.~Redlich, J.~N.~S. Kwong, On the thermodynamics
of solutions. V. An equation of state. Fugacities of gaseous solutions,
Chemical Reviews 44 (1949) 233--244.

\bibitem{rowe1970} A.~M. Rowe Jr., J.~C.~S. Chou, Pressure--volume--
temperature--concentration relation of aqueous NaCl solutions, Journal
of Chemical and Engineering Data 15 (1970) 61--66.

\bibitem{batzle1992} M.~Batzle, Z.~Wang, Seismic properties of pore
fluids, Geophysics 57 (1992) 1396--1408.

\bibitem{garcia2001} J.~E. Garcia, Density of aqueous solutions of
CO$_2$, Report LBNL-49023, Lawrence Berkeley National Laboratory, 2001.

\bibitem{islam2012} A.~W. Islam, E.~S. Carlson, Viscosity models and
effects of dissolved CO$_2$, Energy \& Fuels 26 (2012) 5330--5336.

\bibitem{fenghour1998} A.~Fenghour, W.~A. Wakeham, V.~Vesovic, The
viscosity of carbon dioxide, Journal of Physical and Chemical Reference
Data 27 (1998) 31--44.

\bibitem{davidson1993} T.~A. Davidson, A simple and accurate method for
calculating viscosity of gaseous mixtures, Report of Investigations
9456, US Bureau of Mines, 1993.

\bibitem{juanes2006} R.~Juanes, E.~J. Spiteri, F.~M. Orr Jr.,
M.~J. Blunt, Impact of relative permeability hysteresis on geological
CO$_2$ storage, Water Resources Research 42 (2006) W12418.

\bibitem{killough1976} J.~E. Killough, Reservoir simulation with
history-dependent saturation functions, SPE Journal 16 (1976) 37--48.

\bibitem{land1968} C.~S. Land, Calculation of imbibition relative
permeability for two- and three-phase flow from rock properties, SPE
Journal 8 (1968) 149--156.

\bibitem{bear1972} J.~Bear, Dynamics of Fluids in Porous Media, Dover,
New York, 1972.

\bibitem{flemisch2024} B.~Flemisch, et al., The FluidFlower validation
benchmark study for the storage of CO$_2$, Transport in Porous Media 151
(2024) 865--912.

\bibitem{peaceman1978} D.~W. Peaceman, Interpretation of well-block
pressures in numerical reservoir simulation, SPE Journal 18 (1978)
183--194.

\bibitem{peaceman1983} D.~W. Peaceman, Interpretation of well-block
pressures in numerical reservoir simulation with nonsquare grid blocks
and anisotropic permeability, SPE Journal 23 (1983) 531--543.

\bibitem{palagi1994a} C.~L. Palagi, K.~Aziz, Use of Voronoi grid in
reservoir simulation, SPE Advanced Technology Series 2 (2) (1994) 69--77.

\bibitem{palagi1994b} C.~L. Palagi, K.~Aziz, Modeling vertical and
horizontal wells with Voronoi grid, SPE Reservoir Engineering 9 (1)
(1994) 15--21.

\end{thebibliography}

\end{document}
